\section{Python: lists, tuples and arrays}

It is often convenient to refer to several (or many) pieces of data
via single variable name, e.g., a list of atmospheric
temperatures~\texttt{T} at corresponding heights~\texttt{z}. We can do
this in Python using \textit{lists}, \textit{tuples} or
\textit{arrays}. Lists and tuples are data structures available in
Python whereas arrays are handled using the \texttt{numpy} module.

We can create a (small) list by enclosing the entries in square brackets
and separating them by commas.

\marginnote{
\texttt{>> R = [3, 6, -2]}
}

If we want to access an individual element of a list, either to
compute with it or to assign to it, we can do so using
\textit{indices} (also called \textit{subscripts}) in square brackets
after the list name.  In Python, indices start at~0 (in some
languages they start at~1). We can use negative indices to
access the list (or tuple, or array) from the end.

\marginnote{
\texttt{>> R[0] = 37}
}

Note that square brackets are used here because round brackets are
used for calling functions (you have seen this with built-in functions
e.g. \texttt{sin(pi)} and will meet it again later when you write your
own functions).

\ex{What error message do you get if you try using round
brackets to access the first entry in \texttt{R}? How about if you try
to change the value of the first entry using round brackets?}

A list can be extended by adding a new value at the end using the
\texttt{append} method. Here round brackets are used because we are
calling a function (a method is a function defined on an object).

\marginnote{
\texttt{>> R.append(119)}
}

A Python tuple is very similar to a Python list, but is created using round brackets.

\marginnote{
\texttt{>> R = (3, 6, -2)}
}

Again we can access the individual elements using indices in square brackets.

\marginnote{
\texttt{>> print(R[0])}
}

However, when we try to append to the tuple, or change one of the
values, we find that we cannot. This is because tuples are
\textit{immutable}, meaning that once defined, they cannot be changed.

\ex{Try appending to a tuple, or changing one of the values as
we did with the list above. What error messages do you get?}

The tuples and lists we have just seen look like 1D vectors. However,
if we try to treat them in this way, we will see that they do not
behave as we require. For example, we expect that multiplying a vector
by a scalar will result in a new vector with the same number of
components where each component has been multiplied by the scalar:

\begin{equation}
  a\begin{pmatrix}
  1 \\
  4 \\
  3
  \end{pmatrix}
  = \begin{pmatrix}
  a \\
  4a \\
  3a
  \end{pmatrix}
\end{equation}

If we try to do this with a Python list (or tuple) the result is a new
list (or tuple) that contains $a$ copies of the original list (or
tuple) [Exercise: try it!] We need a different type of object for
performing these mathematical operations... arrays!

The Python module for handling arrays is \texttt{numpy}. The most
common way to import it is to give it the shortened name \texttt{np}.

\marginnote{
\texttt{>> import numpy as np}
}

We can now convert a list to an array

\marginnote{
  \texttt{>> C = np.array([3, 6, -2])}
}

Multiplying this array by a scalar now works as we require for
mathematical operations [Exercise: try it!].

We can create a (small) 2D array, or matrix, in a similar way, with a
new list for each row.

\marginnote{
\texttt{>> M = np.array([[0, 1],  [-1, 0]])}
}

What's the difference between a list, a tuple, a vector, a matrix, and
an array? We usually refer to objects as `vectors' or `matrices' when
we envisage performing matrix-vector arithmetic operations on them
(see below), and as `arrays' when they are just tables of numbers. A
row vector is a $1 \times n$ array with $n$~columns, a column vector
is an $m \times 1$ array with $m$~rows. Python does not distinguish
between the two - they are both represented as 1D arrays. A matrix is
an $m \times n$ array.  Higher dimensional arrays are also
possible. But note all arrays must be `rectangular', i.e.\ have the
same number of columns in every row and vice versa. Note that for 2D
arrays the first index is the row index and the second is the column
index.

\marginnote{
\texttt{>> C[1] = 12}
}

Warning: Try setting \texttt{C[1] = sqrt(3)}. If you print out the
value of \texttt{C[1]} you will see that it is \texttt{1}, rather than
\texttt{1.73205081...}. This is because \texttt{C} was created with
integer entries and Python has assumed that you would like the
entries to remain integers. You can check the type of array entries by
looking at the value of \texttt{C.dtype}, and you can set the type of
entry when you create the array. Try setting \texttt{C = np.array([3,
    6, 2], dtype=float)} and then setting \texttt{C[1] = sqrt(3)}.


We can create larger arrays in various ways. To create an array
containing entries from 1 to 100 with an increment of 1 we can use
\texttt{np.arange}.

\marginnote{
\texttt{>> z = np.arange(1, 100, 1)}
}

The first two numbers passed to \texttt{np.arange} are the start and
the end of the sequence. The final number is the increment, and this
defaults to~1 if left unspecified.

\marginnote{
\texttt{>> z = np.arange(1, 100)}
}

\noindent
Negative and non-integer increments can also be used. If the increment
means that the final value will be missed, then the array terminates
before the number specified by the ending value. For example,
\texttt{np.arange(1, 10, 2)} generates the array \texttt{[1, 3, 5, 7,
    9]}.

A common requirement is for an array to contain a guaranteed start and
end point with evenly spaced elements in between. For this we can use
\texttt{np.linspace(start, end, n)}, where \texttt{n} is the number of
entries.

\marginnote{
\texttt{>> z = np.linspace(1, 5, 10)}
}

\noindent
We can create arrays of zeros or ones using built-in \texttt{numpy} functions.

\marginnote{
\texttt{>> X0 = np.zeros((100,100))}
}

We can perform standard arithmetic operations on arrays provided their shapes
\textbf{conform}, including ...

\marginnote{
\texttt{>> A = np.array([[1, 2], [3, 4]])}
}

\marginnote{
\texttt{>> B = np.array([[5, 6], [7, 8]])}
}

\noindent
addition and subtraction,

\marginnote{
\texttt{>> S = A + B}
}

\noindent
matrix multiplication,

\marginnote{
\texttt{>> P = A @ B}
}

\noindent
element-by-element multiplication,

\marginnote{
\texttt{>> Q = A * B}
}

\noindent
raising each element to a power,
\marginnote{
\texttt{>> Q = A**2}
}

\noindent
and transpose.
\marginnote{
\texttt{>> T = A.T}
}

\vspace{5mm} Note that by `conform' we mean that the matrix dimensions
must be compatible, so for example we can add two $2\times 3$
matrices, and we can multiply a $2\times 3$ matrix by a $3\times 4$
matrix, but not a $2\times 3$ by a $4\times 2$ matrix. Python will
flag an error if the dimensions make the operation impossible. We can
check the shape of a matrix by printing its \texttt{shape} attribute.

\marginnote{
\texttt{>> print(M.shape)}
}

Warning: because \texttt{numpy} arrays are not strictly matrices,
some operations are allowed that you might not expect. When operating
on arrays, \texttt{numpy} compares their shapes to decide if they are
\textit{compatible}. The dimensions are compatible if they are equal
or is one of them is~1. This is best illustrated with an example.

\marginnote{
\texttt{>> X = np.array([[4, 2, 3], [1, 1, 0]])}
}
\marginnote{
\texttt{>> Y = np.array([[2], [1]])}
}
\marginnote{
\texttt{>> X + Y}
}

\texttt{X} has shape \texttt{(2, 3)} and \texttt{Y} has shape
\texttt{(2, 1)} so both dimensions are compatible. In this case,
\texttt{Y} will act like a \texttt{(2, 3)} array and Python will
\textit{broadcast} \texttt{Y[:,0]} to the other columns. The operation
\texttt{X+Y} will add \texttt{Y} to each column of \texttt{X}. You can
read more about this
\href{https://numpy.org/doc/stable/user/basics.broadcasting.html}{here},
although it is outside the scope of this module.

For complex matrices, that is matrices with complex entries, 
the transpose operation \texttt{A.T} becomes a `transpose and take the complex
conjugate of all the  entries' operation, or a so-called \textbf{Hermitian 
transpose}.  This is the `natural' way to transpose a complex matrix, 
as you will find in later courses on algebra.

Many additional linear algebra functions are provided by the
\texttt{numpy.linalg} library. For square matrices \texttt{A}, the
\texttt{linalg.matrix\_power(A, n)} function will raise the matrix
\texttt{A} to the power of \texttt{n}.

\marginnote{
\texttt{>> Y = np.linalg.matrix\_power(A, 2)}
}

\ex{given the arrays \texttt{a}, \texttt{b}, and \texttt{c},}

\marginnote{
\texttt{>> a = np.array([[1, 0, 1], [0, 1, 0]])}
}

\marginnote{
\texttt{>> b = np.array([[6, 5, 4], [3, 2, 1]])}
}

\marginnote{
\texttt{>> c = np.array([-1, 1])}
}

\noindent
which of the following commands are allowed?
\begin{center}
\begin{tabular}{p{30mm} p{30mm} p{30mm}}
\texttt{>> a + b}     & \texttt{>> a + c}     & \texttt{>> a + b.T} \\
\texttt{>> a @ b}       & \texttt{>> b @ a}       & \texttt{>> b @ c}    \\
\texttt{>> b @ a.T}      & \texttt{>> c @ b}      & \texttt{>> a * b}   \\
\end{tabular}
\end{center}
Try to work it out before checking using Python.

Many built-in functions such as \texttt{cos}, \texttt{exp}, etc., can be applied
to arrays. They act element-by-element.

We can perform calculations on sub-arrays using the colon notation to
refer to a range of subscripts. These are often referred to as
\textbf{slices} of an array, with analogies to serving cake.

\marginnote{
\texttt{>> A = np.array([[1, 2, 3], [4, 5, 6]])}
}

\marginnote{
\texttt{>> A[0, :2] = 0}
}

\marginnote{
\texttt{>> A[:, 0] = -1}
}

Note that the notation \texttt{[:2]} refers to all entries up to but \textit{not including} the entry with index~2.

\section{Python: reading and saving data}
\label{sec:readandsave}

We can save \texttt{numpy} arrays to text files using the
\texttt{np.savetxt} command

\marginnote{
\texttt{>> X = np.array([[4, 2, 3], [6, 7, 8]])}
}
\marginnote{
\texttt{>> np.savetxt('my\_array.txt', X)}
}

\vspace{5mm}
\noindent
Python can read this data back in using the \texttt{np.loadtxt}
command.

\marginnote{
\texttt{>> Y = np.loadtxt('my\_array.txt')}
}

\noindent
The \texttt{savetxt} and \texttt{loadtxt} commands have many options; use
\texttt{help} if you want to find out more.


\section{Python: plotting}

Python has extensive graphics capability via the module
\texttt{matplotlib}. Here we will just use a couple of examples to
introduce some of the basics.

First we need to import the functions we need from the module. The
standard way to do this is with the command

\marginnote{
  \texttt{import matplotlib.pyplot as plt}.
}

\eg{A particle's position $x$ is given as a function of time $t$ as
$x = 2.5 \cos(t)$. Plot its position vs time for $t \in [0,10]$.}

\noindent
First create a uniformly spaced list of times between $0$ and $10$: 
this command line 
specifies $101$ points from $0$ to $10$ going \texttt{[0, 0.1, 0.2, $\ldots$, 9.9, 10]}: 

\marginnote{
\texttt{>> t = np.linspace(0, 10, 101)}
}

\noindent
Compute the corresponding positions; note that we need to use the
\texttt{cos} function defined in \texttt{numpy} to apply the
\texttt{cos} function element-by-element to compute a whole array of
$x$'s; the \texttt{cos} function defined in the \texttt{math} module
does not know about arrays.

\marginnote{
\texttt{>> x = 2.5*np.cos(t)}
}

\noindent
Plot the data.

\marginnote{
\texttt{>> plt.plot(t, x)}
}

\noindent
The plot has been created but before we display it, let us add axis
labels and a title.

\marginnote{
\texttt{>> plt.xlabel('t')}
}

\marginnote{
\texttt{>> plt.ylabel('x')}
}

\marginnote{
\texttt{>> plt.title('Position vs time')}
}

\vspace{8mm}
\noindent
Now display the plot. You should see something like
figure~\ref{fig:basicplot}.

\marginnote{
  \texttt{>> plt.show()}
}

\begin{figure}
\hspace{-1cm}
\includegraphics[width=8cm]{basicplot_xvst.png}
\caption{\label{fig:basicplot}
A basic plot.
}
\end{figure}



\eg{Download the file \texttt{USSA.txt} from the ELE page to the current
folder. Let us read the data set}

\marginnote{
\texttt{>> data = np.loadtxt('USSA.txt')}
}

\noindent
and unpack the data; first column of the array contains heights,
second column contains temperatures. Convert the heights from
$\mathrm{m}$ to $\mathrm{km}$.

\marginnote{
\texttt{>> z = data[:, 0]/1000}
}

\marginnote{
\texttt{>> T = data[:, 1]}
}

\vspace{5mm}

\noindent
Plot the data. Indicate data points by black circles; join with black dotted
lines. Make the line a bit thicker for clarity.

\marginnote{
\texttt{>> plt.plot(T, z, 'ko:', linewidth=1.5)}
}

\noindent
Add axis labels and a title.

\marginnote{
\texttt{>> plt.xlabel('T (K)')}
}

\marginnote{
\texttt{>> plt.ylabel('z (km)')}
}

\marginnote{
\texttt{>> plt.title('US Standard Atmosphere')}
}

\vspace{8mm}
\noindent
Now display the plot.

\marginnote{
\texttt{>> plt.show()}
}

\vspace{5mm}
\noindent
Now let us linearly interpolate to a uniformly spaced set of points between
the fourth and fifth data points in the set. Note that here (contrast
section~\ref{sec:linear_interp}) \texttt{zhat} and \texttt{That} are arrays.

\marginnote{
\texttt{>> zhat = np.linspace(z[4], z[5], 11)}
}

\marginnote{
\texttt{>> That = ((zhat - z[4])*T[5] }
}

\marginnote{
\texttt{\ \ \ \ \ \ \ \ + (z[5] - zhat)*T[4])}
}

\marginnote{
\texttt{\ \ \ \ \ \ \ \ /(z[5] - z[4])}
}

\vspace{5mm}

\noindent
Now plot the results on top of the previous plot (you will have to
reproduce the previous plot).

\marginnote{
\texttt{>> plt.plot(That, zhat, 'b+')}
}

Python allows you to save plots so that they can be imported into
documents, presentations, posters, etc. On the plot window, click on
the disk icon and you will be able to choose a filename and file
format.

Some care is needed with the choice of file format. JPEG (\texttt{.jpeg} or \texttt{.jpg})
might be familiar because it is the usual format for saving digital images, e.g.\ as taken
by your 'phone camera.
\marginnote{JPEG: Joint Photographic Experts Group.}
However, it is not ideal for capturing sharp, clear line plots and can often result in
fuzzy looking lines. \texttt{.png} is the default option but can also be fuzzy. \texttt{.pdf} is usually better,
though not all formats are compatible with different downstream applications.
\marginnote{There is some interesting mathematics behind these different file formats;
for example, JPEG is based on a \textbf{Discrete Cosine Transform}.
Look them up if you are interested!}

Also note that you can change the appearance of your saved plot by resizing the plot window
before saving. You will notice that the font size for any text annotations does not change
as you resize the plot window, so resizing the window changes the \textit{relative} size of
plot and text. Also note that the saved plot will not look exactly like the plot on the screen;
Python \textit{renders} the plot again for saving. The best advice is to carefully check any saved plot to make sure it is clear and legible
and looks exactly how you want it.


\section{Other interpolation formulae}

\subsection{Cubic interpolation using Lagrange polynomials} 

We have met linear interpolation, which involves fitting a polynomial of degree~1 (a straight line)
through~2 data points and we have seen that the error involved is $O(h^2)$, that is quadratic 
as a function of the distance $h$ between the two data points. 

If we have more data points, we can clearly try to obtain a better fit and so less of an error. In 
general given $p$ data points, we can think of fitting a polynomial of degree $p-1$, namely 
%
\begin{equation}
y = a_{p-1} x^{p-1} + a_{p-2} x^{p-2} + \cdots + a_1 x + a_0, 
\end{equation}
%
through these points; this involves $p$ unknown coefficients $a_0, a_1, \ldots, a_{p-1}$ and we obtain $p$ equations
to fix these, one per data point. 

There are two issues then, namely how do we do this in practice, and what is the error of the numerical
method as a function of the maximum distance $h$ between the data points? 

To do this in practice we can look back at the equations (\ref{eq:lag1}, \ref{eq:lag2}) and try to generalise the way 
we constructed the functions $L_1$ and $L_2$ that we used for linear interpolation, $p=2$. We'll focus on an example with 
$p=4$, namely \emph{cubic interpolation}. We will put a cubic polynomial through 4 evenly spaced data points, and explicitly let us look at the 
points
%
\begin{equation}
(x_1,y_1) =(1,1), \quad 
(x_2,y_2) =(2,2), \quad
(x_3,y_3) =(3,1), \quad
(x_4,y_4) =(4,4). 
\label{eq:lagdata} 
\end{equation}
%
\marginnote{The method goes through in the same way if the data points are not evenly spaced.} %
Let us define the cubic polynomial $\mathcal{L}_1 (x) $ by 
%
\begin{equation}
\mathcal{L}_1 (x) = \frac{(x-x_2)(x-x_3)(x-x_4)}{ (x_1-x_2)(x_1-x_3)(x_1-x_4)} =  \frac{(x-2)(x-3)(x-4)}{ (1-2)(1-3)(1-4)}\, . 
\notag
\end{equation}
%
If we plug $x=x_1$, $x = x_2$, etc., into this we find that
%
\begin{equation}
\mathcal{L}_1 (x_1) =1, \quad
\mathcal{L}_1 (x_2) =0, \quad
\mathcal{L}_1 (x_3) =0, \quad
\mathcal{L}_1 (x_4) =0. 
\notag
\end{equation}
%
so this \emph{Lagrange polynomial} is switched on --- gives $1$ ---- at the first data point and is switched off --- gives $0$ --- at the rest.
Likewise we can define 
%
\begin{align}
 &  \mathcal{L}_2 (x) = \frac{(x-x_1)(x-x_3)(x-x_4)}{ (x_2-x_1)(x_2-x_3)(x_2-x_4)} =  \frac{(x-1)(x-3)(x-4)}{ (2-1)(2-3)(2-4)}\, , \notag \\
 &  \mathcal{L}_3 (x) = \frac{(x-x_1)(x-x_2)(x-x_4)}{ (x_3-x_1)(x_3-x_2)(x_3-x_4)} =  \frac{(x-1)(x-2)(x-4)}{ (3-1)(3-2)(3-4)}\, ,  \notag\\
 &  \mathcal{L}_4 (x) = \frac{(x-x_1)(x-x_2)(x-x_3)}{ (x_4-x_1)(x_4-x_2)(x_4-x_3)} =  \frac{(x-1)(x-2)(x-3)}{ (4-1)(4-2)(4-3)}\, ,  \notag
\end{align}
%
\marginnote{These expressions look terrible but if you look carefully you can see how they are constructed and the idea behind them}%
Each Lagrange polynomial $\mathcal{L}_k$ is one for $x = x_k$ and zero at the other data points. Each of these is a cubic and they are the building 
blocks that we need to construct a cubic curve through the points in (\ref{eq:lagdata}). 

To do this we now use the $y$ values to define
%
\begin{equation}
y(x) = y_1 \, \mathcal{L}_1 (x)  + y_2 \, \mathcal{L}_2 (x)  + y_3 \, \mathcal{L}_3 (x)  + y_4 \, \mathcal{L}_4 (x)  .\notag
\end{equation}
%
This is a mess, but let us write it down to be clear:
%
\begin{align}
y(x) & = 1 \times  \mathcal{L}_1 (x)  + 2\times  \mathcal{L}_2 (x)  + 1\times \mathcal{L}_3 (x)  + 4 \times \mathcal{L}_4 (x) \notag \\
& = 
1 \times \frac{(x-2)(x-3)(x-4)}{ (1-2)(1-3)(1-4)}
+ 2\times \frac{(x-1)(x-3)(x-4)}{ (2-1)(2-3)(2-4)} \notag\\
& + 1 \times  \frac{(x-1)(x-2)(x-4)}{ (3-1)(3-2)(3-4)}  
+ 4 \times\frac{(x-1)(x-2)(x-3)}{ (4-1)(4-2)(4-3)}\, .
\label{eq:yxlagrange}
\end{align}
%
So, let us look at this! First of all, it consists of four cubic polynomials added together, so $y(x)$ is a cubic polynomial. Secondly if we plug in 
say, $x=x_1=1$, by construction the latter three terms, from $\mathcal{L}_2(1), \mathcal{L}_3(1), \mathcal{L}_4(1)$, are all zero, whereas $\mathcal{L}_1(1)=1$ switches on, giving
$y(1) = 1 \times  \mathcal{L}_1(1) = 1$: the curve goes through $(1,1)$. Likewise if we put in $x=x_2=2$ the second term from  $\mathcal{L}_2$ 
switches on while the first, third and fourth are zero, to give $y(2) = 2 \times \mathcal{L}_2(2) = 2$ and the curve goes through $(2,2)$. Likewise the curves passes through $(3,1)$ and $(4,4)$. Using the Lagrange polynomials as building blocks, we have constructed the interpolating cubic polynomial explicitly as $y(x)$ above. 

Now you may be tempted to multiply out the expressions in (\ref{eq:yxlagrange}), which is a very messy process, but one that would result in the 
answer
%
\begin{equation}
 y(x)  = x^3 - 7 x^2 + 15 x - 8. 
\end{equation}
%
\marginnote{You can check this polynomial goes through the data points, and in fact this is how the example was constructed.}%
%
This may seem a bit depressing: we have the answer, but it would have taken ages to multiply out accurately, and (even) your lecturers would have a high
chance of making slips and getting the wrong answer. And, this is just a simple example not a realistic problem. So what have we gained? Well, the beauty is that we have the 
tool of Python (or other computer languages such as Matlab and Julia), and we can use these to set up the Lagrange polynomials above, and to work out the approximation. We 
never need to multiply anything out ourselves, just use Python to evaluate quantities, as we shall see below. In this way, the MTH1003 module aims to combine theory and numerics to give
you, the students, a powerful set of tools for advancing your study of mathematics at university and beyond. 

\subsection{General theory}

Having seen how the example of cubic interpolation works, let us briefly describe the general 
case where we seek the unique polynomial of degree~$p-1$ that passes through $p$ data points
$(x_1, y_1)\ \ldots \ (x_{p}, y_{p})$. Given these data points we define the Lagrange polynomials
$\mathcal{L}_k (x)$ as 
%
\begin{equation}
\mathcal{L}_k (x) = \prod_{j=1,\ j \ne k}^p \, \, \frac{x - x_j}{x_k - x_j}\, . 
\end{equation}
%
\marginnote{The way these general formulae are written may look quite confusing at first, but hopefully if you refer back to the example above you will see how it works. We are familiar with using capital sigma notation, $\sum a_i $, to denote a sum of terms $a_i$, say, and similarly the capital pi notation, $\prod a_i$, is used to  denote the product of terms $a_i$.}%
Here the product $\prod$ is taken over values of $j$ from $1$ to $p$ inclusive, but omitting the value $k$. If you set $p=4$ and try $k=1$ you will
see that the formula gives $\mathcal{L}_1$ as defined above, for example. These Lagrange polynomials have the required property that 
$\mathcal{L}_k(x_k)=1$ and $\mathcal{L}_k(x_j)=0$ for $j\neq k$. They switch off and on in the correct fashion. 

Then the polynomial that passes through all the data points, the interpolating polynomial, is constructed as the sum of the Lagrange polynomials 
weighted by the $y$-values as 
%
\begin{equation}
\label{eq:lagrange_interp}
y = \sum_{k=1}^{p} y_k\,  \mathcal{L}_k (x) .
\end{equation}

\ex{Check that this formula reduces to \eqref{eq:linear_interp} when
$p=2$.}

\ex{Check that \eqref{eq:lagrange_interp} does indeed give a polynomial
of degree $p-1$ and that $y = y_k$ at $x = x_k$ for $k = 1, \ldots, p$.}

\vspace{5mm}

\subsection{Convergence of cubic interpolation} 

Let's repeat the convergence test of section~\ref{sec:linear_interp}. Suppose we have a constant $h>0$ and use
the points
%
\begin{equation}
x_1 = 1 - 4h/3, \quad
x_2 = 1 - h/3, \quad
x_3 = 1 + 2h/3, \quad
x_4 = 1 + 5h/3, \notag
\end{equation}
%
to do cubic Lagrange interpolation. We will compare the interpolated value of
$\sin(\hat{x})$ to the exact value for $\hat{x} = 1$. With these values of $x_j$ and $\hat{x}$
\eqref{eq:lagrange_interp} becomes (eventually)
%
\begin{eqnarray}
\hat{y} & = & - \frac{5}{81} \sin(1 - 4h/3) + \frac{60}{81} \sin(1 - h/3)
\nonumber 
         + \frac{30}{81} \sin(1 + 2h/3) - \frac{4}{81} \sin(1 + 5h/3) .
\nonumber
\end{eqnarray}

\noindent
In Python


\noindent
\texttt{>> c1 = -5/81}

\noindent
\texttt{>> c2 = 60/81}

\noindent
\texttt{>> c3 = 30/81}

\noindent
\texttt{>> c4 = -4/81}

\noindent
\texttt{>> h = 1}

\noindent
\texttt{>> yerr = abs( c1*sin(1 - 4*h/3) + c2*sin(1 - h/3) + \textbackslash}

\noindent
\texttt{\ \ \ \ \ \ \ \ \ c3*sin(1 + 2*h/3) + c4*sin(1 + 5*h/3) - sin(1) )}

\noindent
\texttt{>> h = h/2}%
\marginnote[-15mm]
{Where do these coefficients come from? Well we have 
$$
y(x) = y_1L_1(x) + y_2 L_2(x) + y_3 L_3(x) + y_4 L_4(x)
$$
for the general interpolating cubic, and we set $x = \hat{x} = 1$ and then $y(\hat{x}) = \hat{y}$ is our approximation to $\sin 1$. Thus  
$$
\hat{y} = y_1L_1(1) + y_2 L_2(1) + y_3 L_3(1) + y_4 L_4(1).
$$  
Now here $y_1 = \sin x_1 = \sin (1 - 4h/3)$ and similarly for $y_2$, $y_3$ and $y_4$. The coefficients are given by the appropriate Lagrange polynomials \emph{evaluated at} $x = \hat{x}=1$, so that for example
$$
L_1(1)  = \frac{ (1- x_2)(1-x_3)(1-x_4) } { (x_1-x_2)(x_1-x_3) (x_1-x_4)} = \cdots =  - \frac{5}{81}, 
$$
obtained by putting in $x_1 = 1 - 4h/3$ and similarly $x_2$, $x_3$, $x_4$. 
} 



\noindent
\texttt{>> yerr = abs( c1*sin(1 - 4*h/3) + c2*sin(1 - h/3) + \textbackslash}

\noindent
\texttt{\ \ \ \ \ \ \ \ \ c3*sin(1 + 2*h/3) + c4*sin(1 + 5*h/3) - sin(1) )}

\noindent
etc. One can see how the error {\tt yerr} goes down when $h$ is halved each time. %


It can be shown, though we will not prove it here, that Lagrange interpolation
with $p$ data points and a polynomial of degree~$p-1$ has an error that is $O(h^{p})$ as $h \rightarrow 0$
\textit{provided} the function being interpolated is sufficiently smooth. However,
if the function not sufficiently smooth then higher degree interpolation is not
necessarily more accurate, nor does it necessarily converge faster, than lower
degree interpolation. Exercise: try cubic interpolation on the three functions
given at the end of
section~\ref{sec:converge_linear_interp}; does it converge faster than linear
interpolation?

There are several other flavours of interpolation besides Lagrange: look up
cubic spline and Hermite interpolation. Interestingly, given more data, for example
$p=20$ data points, constructing the interpolating 
%
\marginnote[-15mm]
{We are referring here to increasing $p$ for fixed $h$ --- when the behaviour generally gets poorer, the oscillations wilder. If we fix $p$ and reduce $h$ then for small enough $h$ we should see convergence at $O(h^p)$ as discussed above. This is one of these instances where the \emph{order} in which we take limits --- here $h\to0$ and $p\to\infty$ --- is important. }%
%
polynomial of degree $19$ that goes through the data points is actually not the most sensible route!
Even though this polynomial exists and is easy
enough to programme in a language such as Python (with some experience...), in reality the polynomial 
starts to gain oscillations \emph{between} the data points that get wilder and wilder (this is not a 
technical term) as more data points $(x_j,y_j)$ are incorporated. Cubic interpolation is a sensible method but
interestingly increasing the order $p$ ceases to help after a while; other methods such as cubic splines
avoid this undesirable behaviour. 


\section{Python: scripts and functions}

\subsection{Motivation}
\label{sec:motivation_two_problems}

Consider these two problems:

\vspace{3mm}

\noindent
1. For the SIR-model that we met in section~\ref{sec:built-in}, suppose that
we know $I_0$ and $S_0$.
Using \eqref{eq:Imax} it is straightforward to work out $I_\mathrm{max}$ given
$q$. But how could we work out
what value of $q$ would produce a given $I_\mathrm{max}$?

\vspace{5mm}

\noindent
2. A simple model for the temperature $T$ (Kelvin) and concentration $q$ of
water vapour in a saturated air parcel with given total enthalpy $H$ is
\begin{equation}
q = q_0 \exp(20 - 5000/T), \qquad C_\mathrm{p} T + L q = H,
\end{equation}
where $q_0 \approx 0.0038$,
$C_\mathrm{p} \approx 1000 \, \mathrm{J \, kg}^{-1}\mathrm{K}^{-1}$ is the
specific heat capacity at constant pressure and
$L \approx 2.5 \times 10^6 \, \mathrm{J \, kg}^{-1}$ is the specific latent heat of
vaporisation. We can eliminate $q$ from these equations, but how can we solve
for $T$?
\begin{marginfigure}
\begin{center}
\includegraphics[width=2.5cm]{AdaLovelace.jpg}
\end{center}
\caption{\label{fig:AdaLovelace}
Ada Lovelace (1815--1852) worked with Charles Babbage on his Analytical Engine.
She is widely regarded as one of the world's first computer programmers,
and also foresaw the potential for computing machines beyond `mere
number crunching'.}
\end{marginfigure}


\vspace{5mm}

\noindent
Both problems can be solved using \textit{iterative methods}, which we will meet
shortly, in which certain computations are repeated over and over again, each time
(hopefully) producing a more accurate approximation to the solution.

Such calculations become tedious and impractical when done via the command line;
it would be better if we could give a list of commands to Python and ask it to repeat
them. We can do this using \textit{scripts} and \textit{functions}. 



\subsection{Scripts}

A Python script is a file containing a list of Python commands. The
filename must end in the extension \texttt{.py}

\marginnote{
e.g.\ \texttt{myscript.py}
}

\noindent
The commands in the file can be executed, \textit{in order}, by typing
\texttt{python} followed by the name of the script at the command line
prompt.

\marginnote{
\texttt{>> python myscript.py}
}

\noindent
You can also run scripts from within iPython by typing \texttt{run myscript.py}.

There are many text editors or IDEs available for opening and editing
Python scripts. In this course we will be using Jupyter Lab. To open a
new script, launch Jupyter Lab and choose the `Python File' icon in
the `Other' section. The script will automatially be given the name
\texttt{untitled.py}, but you can change this when you save it (either
by selecting \texttt{Save Python File} from the \texttt{File} menu, or
by pressing the shortcut combination \texttt{Ctrl+S}).

\marginnote{
\texttt{\# My first script}
}

\marginnote{
\texttt{import numpy as np}
}


\marginnote{
\texttt{\# List the numbers from 1 to 10}
}

\marginnote{
\texttt{N = np.linspace(1, 10, 10)}
}

\marginnote{
\texttt{print(N)}
}

\marginnote{
\texttt{\# and display their squares}
}

\marginnote{
\texttt{print(N**2)}
}

\marginnote{
\texttt{} % Blank line to create space before footnote!
}

You can include \textit{comments} in your code to help explain it to
others and to remind yourself what it does; it is highly recommended
that you do this! Everything to the right of a \texttt{\#} symbol is
treated as a comment and is not acted on.

If you need to include a long calculation it can be continued over
several lines; ending a line with a backslash,
i.e. one of these: \texttt{\textbackslash}, will cause the following line to be
treated as a continuation of this one. E.g.

\vspace{3mm}

\noindent
\texttt{\# Cubic interpolation of sin(1)}

\noindent
\texttt{from math import sin}

\noindent
\texttt{\# Coefficients}

\noindent
\texttt{c1, c2, c3, c4 = -5/81, 60/81, 30/81, -4/81}

\noindent
\texttt{\# Interval between data points}

\noindent
\texttt{h = 1}

\noindent
\texttt{\# Error in cubic interpolation formula}

\noindent
\texttt{yerr = abs( c1*sin(1 - 4*h/3) + c2*sin(1 - h/3) + \textbackslash }

\noindent
\texttt{\ \ \ \ \ \ \ \ \ \ \ \ c3*sin(1 + 2*h/3) + c4*sin(1 + 5*h/3) - sin(1) )}

\vspace{3mm}

\noindent
It is also possible to include several commands on one line as I have
done here, but generally your code will be clearer and easier to debug
if you do not do this. The \textbackslash ~is not strictly necessary in
the above example as Python permits linebreaks within brackets.

To execute a script it needs to be in the current folder. A very
common cause of errors is when the script you want to execute is not
in the current folder!

Notice that when you execute a script within iPython, the variables
evolve exactly as if you had typed those same commands on the command
line (but much faster!). This means that if there are any errors in
your script, the script will stop running at this point and give an
error message; any lines after the error will not be executed.


\subsection{Functions}

We have seen some examples of built-in functions in Python
(e.g. \texttt{sin} and \texttt{log}). We can also define our own
functions. Similarly, they will have \textit{inputs} and
\textit{outputs}. E.g., a simple function to evaluate the cubic
\begin{displaymath}
f(x) = x^3 -2x + 6
\end{displaymath}
is

\vspace{3mm}

\noindent
\texttt{def cubic(x):}

\noindent
\texttt{\ \ \ \# CUBIC Simple cubic function}

\noindent
\texttt{\ \ \ b = -2}

\noindent
\texttt{\ \ \ c = 6}

\noindent
\texttt{\ \ \ f = x**3 + b*x + c}

\noindent
\texttt{\ \ \ return f}

\vspace{3mm}

\noindent
This function can be saved in the file \texttt{cubic.py}.

Note how the input \texttt{x} is specified in the function
\textit{interface}, and note that you need to be careful to get the
correct use of parentheses \texttt{()} on the right-hand side.  When
the variable \texttt{x} appears on the right hand side of a
calculation within the function, the input value is used there.  The
value computed and stored in the variable \texttt{f} will be output by
the function.

The function is used in the same way as built-in functions. In this example,
2 is passed into the function and used as the value of~\texttt{x}.
The value of $f(2)=2^3-2*2+6$ is passed as the function output and is displayed.

\marginnote{
\texttt{>> y = cubic(2)}
}

\marginnote{
\texttt{>> print(y)}
}

%\vspace{3mm}

\noindent
In the next example,
the value of \texttt{z}, i.e.\ 3.3, is passed into the function and used as
the value of~\texttt{x}. The value of $f$ is passed as the function output
and is stored in the variable~\texttt{w}.

\marginnote{
\texttt{z = 3.3}
}

\marginnote{
\texttt{w = cubic(z)}
}

Function inputs and outputs can be arrays. In fact
this function will work for arrays too, because
\texttt{x**3} will work equally well for a scalar \texttt{x} or an array.
Thus, this function will output an element-by-element evaluation of the
cubic $f(x)$ when its input is an array.

\marginnote{
\texttt{t = np.linspace(-5, 5, 101)}
}

\marginnote{
\texttt{q = cubic(t)}
}

\marginnote{
\texttt{plt.plot(t, q, 'ro')}
}
\marginnote{
\texttt{plt.show()}
}


In this example we chose to call the function input argument
\texttt{x} and the output argument \texttt{f}. However, we could have
used any other variable name for the input; the crucial point is that
the input variable name appearing in the first line of the function
file should agree with that used in the calculation. Similarly, we
could have used any other variable name for the output argument. For
this reason, input and output arguments are sometimes referred to as
\textit{dummy arguments}.

\marginnote{
Make sure you understand how function input and output arguments work.
}

Functions can have multiple inputs and/or multiple outputs. E.g.

\vspace{3mm}

\noindent
\texttt{def argand(rr, ii):}

\noindent
\texttt{\ \ \ \#ARGAND Compute the modulus and argument of a complex number}

\noindent
\texttt{\ \ \ \# with real part rr and imaginary part ii}

\noindent
\texttt{\ \ \ from math import sqrt, atan2};

\noindent
\texttt{\ \ \ mm = sqrt(rr*rr + ii*ii)}

\noindent
\texttt{\ \ \ aa = atan2(ii, rr)}

\noindent
\texttt{\ \ \ return mm, aa}

\marginnote{
Use \texttt{help} to find out about the built-in \texttt{atan2} function.
Note that it has two input arguments.
}

\vspace{3mm}

A crucial point about functions is what happens with the
namespace. When you call a function a new \textit{local} namespace is
created, and all the variables used by the function exist only in this
local namespace. When Python exits the function this local namespace
disappears.

Although this idea may seem complicated at first, in fact it is very powerful.
It means that the inner workings of a function are invisible outside of it
so, for example, there is no danger of a clash of variable names. For this reason
it is often better practice to use functions rather than scripts in your coding.


\subsection{More on scripts and functions}

The allowed names for scripts and functions are exactly the same as those
for variables. This flexibility is useful, but it means there is a danger
of name clashes between variables, your own scripts and functions,
and built-in functions; the onus is on you to check and avoid such clashes.

We have seen that we can run scripts and call functions from the
iPython command line.  In fact scripts can call functions, and
functions can call other functions. In this way we can build up
complex computations from simple components; indeed it is good
practice to break down complex computations into simple components
that can be tested and debugged individually.


\section{Python: flow control}

So far we have used Python commands, executed in order, to do simple
calculations and to read, write, and plot data. Scripts and functions
become very powerful when we use additional commands to
\textit{control the flow} of the computation, i.e.\ to determine which
commands are executed and allowing repetition of commands.

\subsection{\texttt{for} loops}

\texttt{for} loops allow us to repeat a section of code a number of times. The generic
syntax is

\vspace{3mm}

\noindent
\texttt{for} \textit{<variable>} \texttt{=} \textit{<expression>}:

\noindent
\texttt{\ \ \ ...}

\noindent
\textit{\ \ \ commands}

\noindent
\texttt{\ \ \ ...}

\vspace{1mm}

\marginnote{
E.g.\
}

\marginnote{
\texttt{for k in range(10):}
}

\marginnote{
\texttt{\ \ \ \ print(k, k**3)}
}

\vspace{3mm}

\noindent
The \textit{<expression>} should evaluate to to a Python object that is \textit{iterable}, i.e. it can be iterated over, like a list.
The \textit{<variable>} is set to the first element of the iterable
and the \textit{commands} are executed. The \textit{<variable>} is then set to the
second element of the iterable and the \textit{commands} are executed, and so on
until all elements of the iterable have been used.
The \textit{<expression>} may evaluate to an empty iterable (with zero elements),
in which case the loop is executed zero times, for example:
 
\marginnote{
\texttt{n=0}
}

\marginnote{
\texttt{for k in range(n):}
}

\marginnote{
\texttt{\ \ \ \ print(k, k**3)}
}


\subsection{while loops}

\texttt{while} loops also allow us to repeat a section of code a number of times.
The generic synatax is

\vspace{3mm}

\noindent
\texttt{while} \textit{<logical expression> :}

\noindent
\texttt{\ \ \ ...}

\noindent
\textit{\ \ \ commands}

\noindent
\texttt{\ \ \ ...}

\marginnote{
E.g.\
}

\marginnote{
\texttt{total = 0}
}

\marginnote{
\texttt{N = 0}
}

\marginnote{
\texttt{while total < 1000:}
}

\marginnote{
\texttt{\ \ \ \ N = N + 1}
}

\marginnote{
\texttt{\ \ \ \ total = total + N**2}
}

\marginnote{
\texttt{print(total)}
}


\vspace{3mm}

\noindent
If the \textit{<logical expression>} evaluates to \texttt{False} then the \textit{commands}
are not executed. If the \textit{<logical expression>} evaluates to \texttt{True}
then the commands are executed; the \textit{<logical expression>} is then
re-evaluated and the process repeats.
Clearly the \textit{commands} must
be able to alter the result of the \textit{<logical expression>}, otherwise
the loop will repeat for ever.

\marginnote{
If you want to interrupt the execution of a Python script or function, e.g.\ 
because it is stuck in an infinite loop, then use \textit{ctrl-c}.
}


\subsection{\texttt{if} blocks}

\texttt{if} blocks allow the code to decide, at run time, whether or not to
execute certain blocks of commands. The simplest version of an \texttt{if}
block has the syntax

\vspace{3mm}

\noindent
\texttt{if} \textit{<logical expression>:}

\noindent
\texttt{\ \ \ ...}

\noindent
\textit{\ \ \ commands}

\noindent
\texttt{\ \ \ ...}

\marginnote{
E.g.\
}

\marginnote{
\texttt{n = 13}
}

\marginnote{
\texttt{if n\%2 == 0:}
}

\marginnote{
\texttt{\ \ \ \ print('n is even')}
}

\vspace{3mm}

\noindent
If the \textit{<logical expression>} evaluates to true
then the commands are executed, otherwise they are not. Note that,
in contrast to a \texttt{while} loop, there is no looping or re-evaluation
of the \textit{<logical expression>} implied by the \texttt{if} block.

There are several more elaborate variants of \texttt{if} blocks, which work
in the obvious way:

\vspace{2mm}

\begin{tabular}{p{4.3cm}p{4.7cm}p{4.5cm}}
\texttt{if} \textit{<logical expression>:} & \texttt{if} \textit{<logical expression>:} &
\texttt{if} \textit{<logical expression>:} \\
\texttt{\ \ \ ...} & \texttt{\ \ \ ...} & \texttt{\ \ \ ...} \\
\textit{\ \ \ commands-1} & \textit{\ \ \ commands-1} & \textit{\ \ \ commands-1} \\
\texttt{\ \ \ ...} & \texttt{\ \ \ ...} & \texttt{\ \ \ ...} \\
\texttt{else:} & \texttt{elseif} \textit{<logical expression>:} &
\texttt{elseif} \textit{<logical expression>:} \\
\texttt{\ \ \ ...} & \texttt{\ \ \ ...} & \texttt{\ \ \ ...} \\
\textit{\ \ \ commands-2} & \textit{\ \ \ commands-2} & \textit{\ \ \ commands-2} \\
\texttt{\ \ \ ...} & \texttt{\ \ \ ...} & \texttt{\ \ \ ...} \\
& \texttt{elseif} \textit{<logical expression>:} & \texttt{else:} \\
             & \texttt{\ \ \ ...} & \texttt{\ \ \ ...} \\
             & \textit{\ \ \ commands-3} & \textit{\ \ \ commands-3} \\
             & \texttt{\ \ \ ...} & \texttt{\ \ \ ...} \\
             & & 
\end{tabular}

\vspace{1mm}


\marginnote{
\noindent
\
}

\marginnote{
\noindent
\
}

\marginnote{
\noindent
E.g.\
}

\marginnote{
\texttt{n = 13}
}

\marginnote{
\texttt{if n\%2 == 0:}
}

\marginnote{
\texttt{\ \ \ \ print('n is even')}
}


\marginnote{
\texttt{else:}
}

\marginnote{
\texttt{\ \ \ \ print('n is odd')}
}


%\newpage

\subsection{More on flow control}

You will notice that in the above examples the code inside loops and
if blocks is indented. This is crucial in Python as there is no
\texttt{end} statement (unlike in some other coding languages). Python
will complain at indentation that it cannot parse, but you need to
take care that you only indent the code that you intend to.

Python allows you to build up complex programmes by putting loops within loops,
loops within if blocks, if blocks within loops, and if blocks within if blocks.
You can use loops and if blocks within scripts and functions, and you can call
scripts and functions from within loops and if blocks.


\section{Algorithms}

It is often tricky to see the connection between a set of equations comprising
a mathematical model and the computer code needed to solve them. Don't start
typing code too soon; first make sure you understand the \textit{algorithm}.

\marginnote{
The word \textit{algorithm} is the latinisation of the name of an ancient scholar from Khw\={a}rezm in Central Asia. See \href{https://earthobservatory.nasa.gov/images/91544/how-algorithm-got-its-name}{this article}.
}

\vspace{3mm}

\noindent
An \textbf{algorithm} is a list of rules or steps to be followed in order to
perform a computation or solve a problem.

\vspace{1mm}


\begin{itemize}
\item
An algorithm must be unambiguous.
\item
The order of the steps matters.
\item
It is often helpful to express an algorithm in terms of what is called \textbf{pseudocode}.
\end{itemize}

\noindent
As an example of pseudocode, figure \ref{fig:pseudocode1} shows pseudocode for the task of listing the legal
moves in a game of noughts and crosses. Figure \ref{fig:pseudocode2} shows pseudocode
for the task of listing the winning moves in a game of noughts and crosses.

%\begin{figure}
%\hspace{-3cm}
%\includegraphics[width=17cm]{flowchart.png}
%\caption{\label{fig:flowchart}
%Flowchart for the task of listing the legal moves in noughts and crosses.
%}
%\end{figure}

\begin{figure}
{\tt

read in the grid

loop over columns

\ \ \ \ loop over rows

\ \ \ \ \ \ \ \ if this cell is empty

\ \ \ \ \ \ \ \ \ \ \ \ add its coordinates to a list of empty cells

\ \ \ \ end loop over rows

end loop over columns

\  

loop over entries in the list 

\ \ \ \ print out coordinates for each entry

end loop over list

}
\caption{\label{fig:pseudocode1}
Pseudocode for the task of listing the legal moves for X
in noughts and crosses.
}
\end{figure}

\noindent
In each case the procedure is not written in a formal computer language, 
but is designed to help fix ideas, sketch out key elements of the code and the flow
control, and to communicate the procedure (to oneself or to other people). It is 
often a useful step in designing a script before typing it in, and is also useful for 
checking/debugging  a script when it fails.

\begin{figure}
{\tt

loop over columns

\ \ \ \ loop over rows

\ \ \ \ \ \ \ \ if legal

\ \ \ \ \ \ \ \ \ \ \ \ if 2 X's in this column $\rightarrow$ win

\ \ \ \ \ \ \ \ \ \ \ \ if 2 X's in this row $\rightarrow$ win

\ \ \ \ \ \ \ \ \ \ \ \ if first diagonal

\ \ \ \ \ \ \ \ \ \ \ \ \ \ \ \ if 2 X's on this diagonal $\rightarrow$ win

\ \ \ \ \ \ \ \ \ \ \ \ if second diagonal

\ \ \ \ \ \ \ \ \ \ \ \ \ \ \ \ if 2 X's on this diagonal $\rightarrow$ win

\ \ \ \ \ \ \ \ \ \ \ \ if win

\ \ \ \ \ \ \ \ \ \ \ \ \ \ \ \ add this move to the list

\ \ \ \ end loop over rows

end loop over columns

}
\caption{\label{fig:pseudocode2}
Pseudocode for the task of listing the winning moves for X
in noughts and crosses.
}
\end{figure}


Algorithms are behind any process designed to carry out a task, make a calculation, or to come to a conclusion. For example one could think 
of a court case involving defendant, prosecution, judge and jury as an algorithm to determine guilt or innocence: there is a 
clear sequence of steps, taken in a careful order. Other algorithms could include making breakfast: get a bowl, put 
in some cereal, add milk to the cereal, eat the food, wash up. Another algorithm could be the sequence of steps a mechanic might take to 
diagnose a car that won't start: check the battery, check the starter motor, ... Algorithms are everywhere! However here of course we will focus on algorithms for computations, tasks linked to mechanics, and mathematical modelling. 

\ex{Look up Euclid's algorithm for finding the greatest common divisor
of two integers. Express the algorithm as pseudocode.}


\section{Solving nonlinear equations}
\label{sec:solve}

Let us return to the two problems that we posed in
section~\ref{sec:motivation_two_problems}. Both problems can be re-expressed
as a requirement to solve
a nonlinear equation of the form $ f(x) = 0 $ for some scalar function $ f $
of a scalar variable $ x $. In this section we will look at some
\textit{iterative methods}, which generate
a sequence of approximate solutions
$ \left\{ x^{(m)} ; \ m=0, 1, 2, ... \right\} $. In order to be useful
a method must \textit{converge} to the
true solution $ x^{\mathrm{*}} $, that is
\begin{displaymath}
x^{(m)} \rightarrow x^{\mathrm{*}} \mathrm{\ \ as\ \ } m \rightarrow \infty. 
\end{displaymath}

\marginnote{
Compare this definition of \textit{convergence} with the one we used in
section~\ref{sec:converge_linear_interp}.
}


\subsection{Fixed point iteration}
\label{sec:fpi}



One of the simplest iterative methods is \textit{fixed point iteration}.
In this method we 
rewrite $ f(x) = 0 $ in the form $ g(x) = x $. (In general there will be many
possible ways to do this; a good choice is to make $g'(x^*)$ small---see below.)
Then we simply iterate:
\begin{displaymath}
x^{(m+1)} = g(x^{(m)}) ,
\end{displaymath}
starting from some first guess $x^{(0)}$. 
If $x^{(m)}$ converges to some value $x^*$ then so does $x^{(m+1)}$ and so 
the limit of the above  equation yields $x^* = g(x^*)$, and we have a solution. That is, ... \textbf{if}  
there is convergence. Possible pseudocode for this process is seen in figure \ref{fig:flowchart_fpi}. 

%\begin{figure}
%\hspace{-2cm}
%\includegraphics[width=14cm]{flowchart_fpi.png}
%\caption{\label{fig:flowchart_fpi}
%Flowchart for fixed point iteration.
%}
%\end{figure}

\begin{figure}
{\tt

set the initial value $x^{(0)}$

loop over number of iterations 

\ \ \ \ \ set $x^{(m)} = g(x^{(m-1)})$ 

\ \ \ \ \ if stopping criterion has been met 

\ \ \ \ \ \ \ \ \ leave the loop

end loop over iterations

print out final answer  $x^{(m)}$

}
\caption{\label{fig:flowchart_fpi}
Pseudocode for fixed point iteration.
}
\end{figure}


In practice we cannot iterate for ever; we must decide when to stop the iteration.
This may be, for example,
\begin{itemize}
\item
after a fixed number of iterations;
\item
when $|g(x^{(m)})|$ is small enough;
\item
when the change in $x$, $|x^{(m+1)} - x^{(m)}|$, is small enough.
\end{itemize}

Once we have decided on a stopping criterion, the algorithm for fixed point
iteration is very straightforward.

\eg{ Let's try to solve (a)~$\cos(x) = x$ using fixed point
iteration. A Python script might look like this.}

\marginnote[-13.5mm]{
\texttt{from math import cos}
}

\marginnote[-10mm]{
\texttt{\# Set first guess}
}

\marginnote[-6.5mm]{
\texttt{oldx = 0}
}

\marginnote[-3mm]{
\texttt{x = 1}
}

\marginnote{
\texttt{\# Iterate until difference is}
}

\marginnote{
\texttt{\# small enough}
}

\marginnote{
\texttt{while abs(x - oldx) > 1e-4:}
}

\marginnote{
\texttt{\ \ \ \ oldx = x}
}

\marginnote{
\texttt{\ \ \ \ x = cos(oldx)}
}

\marginnote{
\texttt{print(f'Root is \{x\}')}
}

\noindent
Notes:
\begin{enumerate}
\item
Although the mathematical description has a superscript $(m)$ on $x^{(m)}$,
it is not necessary to save all of the $x$'s in an array (though we could if
we wanted to).
\item
Here we have used a stopping criterion based on the difference between the
last two $x$'s, so we need to keep track of these in the variables \texttt{oldx}
and \texttt{x}. Notice how we update the values of \texttt{oldx}
and \texttt{x} at each iteration.
\marginnote{
Note that initially we have to set \texttt{oldx} to a value suitably different from \texttt{x}
so that the code within the \texttt{while} is sure to be executed at least once.
}
\item
The algorithm here is slightly different from that in the pseudocode because
the \texttt{while} loop checks whether the stopping criterion has been met
\textit{before} computing the new value of $x$.
\end{enumerate}


We can get an idea of how quickly fixed point iteration converges for this problem,
for example by printing out the values of \texttt{x} at each iteration, or by including a
variable that counts the number of iterations. In this example the code stops
after 23~iterations.

\ex{Modify the code to solve the problems (b)~$\cos(x/2) = x$;
  (c)~$\cos(2x) = x$.  What do you notice about the convergence of
  fixed point iteration for these three problems?}

\vspace{3mm}

We can analyse the convergence of the fixed point iteration method as follows.
Let $ x^{\mathrm{*}} $ be the true solution, so that $x^{\mathrm{*}} = g(x^{\mathrm{*}})$.
Let $\varepsilon^{(m)} = x^{(m)} - x^{\mathrm{*}}$ be the error in our estimate after
$m$ iterations, and assume $\varepsilon^{(m)}$ is small.
\marginnote{Here we don't take the absolute value of $\varepsilon^{(m)}$.}
Then the error after one
further iteration is
\begin{eqnarray}
\varepsilon^{(m+1)} & = & x^{(m+1)} - x^{\mathrm{*}} \nonumber \\
                    & = & g(x^{(m)}) - x^{\mathrm{*}} \nonumber \\
                    & = & g( x^{\mathrm{*}} + \varepsilon^{(m)} ) - x^{\mathrm{*}} \nonumber \\
                    & \approx & g( x^{\mathrm{*}} ) + g'(x^{\mathrm{*}}) \varepsilon^{(m)}
                                                  - x^{\mathrm{*}} \nonumber \\
                    & = & g'(x^{\mathrm{*}}) \varepsilon^{(m)} \nonumber
\end{eqnarray}
Note we have used the Taylor series to approximate\footnote{
We say we have `\textit{linearized}' the function $g(x)$.
Linearization is a hugely important idea: 
you will meet it several times in this module and many times
in your mathematical career!}
$g( x^{\mathrm{*}} + \varepsilon^{(m)} )$.
Thus, the error grows or decays by a factor $g'(x^{\mathrm{*}})$ at every iteration.
A necessary condition for convergence is $|g'(x^{\mathrm{*}})| < 1$. Smaller
$g'(x^{\mathrm{*}})$ gives faster convergence.

\ex{Is this theoretical prediction consistent with what you found in
  practice for the three problems (a), (b) and (c) above?}

\eg{For the SIR epidemic model (section~\ref{sec:built-in},
  \ref{sec:motivation_two_problems}), take $I_0 = 10^3$, $S_0 = 6
  \times 10^7$; use fixed point iteration to determine the value of
  $q$ required to keep $I_\mathrm{max}$ down to $10^5$.}

\eg{For the saturated air parcel model
  (section~\ref{sec:motivation_two_problems}), eliminate $q$ from the
  equations and use fixed point iteration to determine the value of
  $T$ that gives $H = 3.61 \, \mathrm{J \, kg}^{-1}$.}

\vspace{3mm}


You should find that you can solve one of these examples with a straightforward
application of fixed point iteration, but the other does not converge.
So what can we do? One option is to rearrange the problem $x = g(x)$ into a
form with a different $g(x)$; there are always infinitely many ways to do
this, but it requires some insight to make a good choice. Alternatively we
can choose a different, more robust, method to solve the problem.


\subsection{Bracketing methods}

Suppose we wish to solve $ f(x) = 0 $, and we know that $f(x)$ is continuous.


\subsubsection{Bisection}

Suppose we have $ x^{(\mathrm{l})} < x^{(\mathrm{r})} $ (superscripts $(\mathrm{l})$
and $(\mathrm{r})$ indicate left and right)
with $ f(x^{(\mathrm{l})}) $ and
$ f(x^{(\mathrm{r})}) $ of opposite sign. Continuity of $f$ implies that
there must be a zero between $x^{(\mathrm{l})}$ and $x^{(\mathrm{r})}$,
that is, $x^{(\mathrm{l})}$ and $x^{(\mathrm{r})}$ must \textit{bracket} a root.

Let
\begin{displaymath}
x^{(\mathrm{n})} = \frac{1}{2}\left( x^{(\mathrm{l})} + x^{(\mathrm{r})} \right),
\end{displaymath}


% Start marginnotes here to fit them on the page.

\marginnote{
\texttt{\# Solve myf(x)=0 using bisection}
}

\marginnote{
\texttt{\# First guess for interval}
}

\marginnote{
\texttt{xl = 0}
}

\marginnote{
\texttt{fl = myf(xl)}
}

  
\marginnote{
\texttt{xr = 1}
}

\marginnote{
\texttt{fr = myf(xr)}
}

%\marginnote{
%\texttt{\% Check this brackets a root}
%}

%\marginnote{
%\texttt{if xr < xl}
%}

%\marginnote{
%\texttt{\ \ \ \ disp(['xr should be > xl'])}
%}

%\marginnote{
%\texttt{\ \ \ \ return}
%}

%\marginnote{
%\texttt{end}
%}

%\marginnote{
%\texttt{if fl*fr > 0}
%}

%\marginnote{
%\texttt{\ \ \ \ disp(['Choose a different first guess'])}
%}

%\marginnote{
%\texttt{\ \ \ \ return}
%}

%\marginnote{
%\texttt{end}
%}

\marginnote{
\texttt{\# Iterate bisection method}
}

\marginnote{
\texttt{while xr - xl > 1e-4:}
}

\marginnote{
\texttt{\ \ \ \ xn = 0.5*(xl + xr)}
}

\marginnote{
\texttt{\ \ \ \ fn = myf(xn)}
}

\marginnote{
\texttt{\ \ \ \ if fn*fl <= 0:}
}

\marginnote{
\texttt{\ \ \ \ \ \ \ \ \# Root is in left}
}

\marginnote{
\texttt{\ \ \ \ \ \ \ \ \# half interval}
}

\marginnote{
\texttt{\ \ \ \ \ \ \ \ xr = xn}
}

\marginnote{
\texttt{\ \ \ \ \ \ \ \ fr = fn}
}

\marginnote{
\texttt{\ \ \ \ else:}
}

\marginnote{
\texttt{\ \ \ \ \ \ \ \ \# Root is in right}
}

\marginnote{
\texttt{\ \ \ \ \ \ \ \ \# half interval}
}

\marginnote{
\texttt{\ \ \ \ \ \ \ \ xl = xn}
}

\marginnote{
\texttt{\ \ \ \ \ \ \ \ fl = fn}
}

\marginnote{
\texttt{print(f'Root lies in [\{xl\}, \{xr\}]')}
}


\noindent
(superscript $(\mathrm{n})$ for new).
Retain whichever of \( f(x^{(\mathrm{l})}) \) and \( f(x^{(\mathrm{r})}) \) has the
opposite sign to \( f(x^{(\mathrm{n})}) \) and discard the other;
rename the retained points as \( x^{(\mathrm{l})} \), \( x^{(\mathrm{r})} \).
The two retained points still bracket the root and are closer than the original 
two points (half the distance). We now iterate, that is repeat the procedure.

This method is robust; provided $f(x)$ is continuous and the initial values of
$ f(x^{(\mathrm{l})}) $ and $ f(x^{(\mathrm{r})}) $ have opposite sign, the method is
guaranteed to converge to a solution.
However, it is quite slow to converge: the bracket size is halved at each iteration.

A basic Python implementation of the bisection method might look like this.

\ex{Add some code to check that the first guess satisfies $
  x^{(\mathrm{l})} < x^{(\mathrm{r})} $ and $ f(x^{(\mathrm{l})})
  f(x^{(\mathrm{r})}) \le 0$.}

\subsubsection{Secant method}

Suppose we have \( x^{(m-1)} \), \( x^{(m)} \) and the corresponding
 \( f(x^{(m-1)}) \) and \( f(x^{(m)}) \). 
(Here the superscript $(m)$ counts the iterations).
Let
\begin{displaymath}
x^{(m+1)} = \frac{x^{(m-1)}f(x^{(m)}) - x^{(m)}f(x^{(m-1)})}{f(x^{(m)}) - f(x^{(m-1)})}.
\end{displaymath}
Here $x^{(m+1)}$ is an estimate for the zero of $f$ obtained by fitting a
straight line through the pair of points $(x^{(m-1)} , f(x^{(m-1)}))$,
$(x^{(m)} , f(x^{(m)}))$. \marginnote{Exercise: check the secant method formula
for $x^{(m+1)}$; see if you can derive it from the
equation for linear interpolation \eqref{eq:linear_interp} with $\hat{y} = 0$.}
Then discard the oldest data point \( x^{(m-1)} \), \( f(x^{(m-1)}) \) and iterate.

This usually converges much quicker than the bisection method,
but can behave badly if $ f(x^{(m)}) - f(x^{(m-1)}) $ becomes small.
It is not really a bracketing method as the root is not guaranteed
to lie between $ x^{(m)} $ and $ x^{(m+1)} $.


\subsubsection{Regula falsi}

Like the bisection method, regula falsi is a bracketing method.
However, the new value $ x^{(\mathrm{n})} $ is constructed as in the
secant method:
\begin{displaymath}
x^{(\mathrm{n})} = \frac{x^{(\mathrm{l})}f(x^{(\mathrm{r})})
- x^{(\mathrm{r})}f(x^{(\mathrm{l})})}{f(x^{(\mathrm{r})}) - f(x^{(\mathrm{l})})}.
\end{displaymath}
Again, we retain whichever of $ f(x^{(\mathrm{l})}) $ and $ f(x^{(\mathrm{r})}) $ has the
opposite sign to $ f(x^{(\mathrm{n})}) $ and discard the other;
rename as $ x^{(\mathrm{l})} $, $ x^{(\mathrm{r})} $, etc., and iterate.
This ensures that
the two retained values of $ x $ always bracket the root, making this
method more robust than the secant method. However, convergence
is usually somewhat slower.

\ex{Modify the bisection method Python code given above to use the
  regula falsi method instead.}

\ex{You should now be able to solve both the SIR model problem and the
  saturated air parcel problem posed in section~\ref{sec:fpi}, if not
  with fixed point iteration then with one of the bracketing methods.}

